%%
%% Test: Theorem Environments — v1.1
%% Stage: theorem refactor (stages 1-4)
%% 
%% Purpose: Verify theorem system with amsthm + tcolorbox:
%%   - All environments work (theorem, lemma, corollary, proposition)
%%   - Definition, example, remark, exercise
%%   - Solution (no numbering)
%%   - Proof with \qed symbol
%%   - Cleveref references work
%%   - Shared counter (theorem family)
%% 

\documentclass{matinbook}

\title{تست محیط‌های قضیه — نسخه ۱.۱}
\author{تیم MatinBook}
\date{۱۴۰۵}

\booktitle{تست قضیه‌ها}

\begin{document}

\maketitle

\chapter{بررسی محیط‌های قضیه}
\label{chap:test-theorem}

این سند، محیط‌های قضیه \texttt{mb-theorem} را بررسی می‌کند.

%----------------------------------------
\section{خانواده قضیه (شمارنده مشترک)}
\label{sec:theorem-family}

\begin{definition}[عدد اول]
    عدد طبیعی بزرگ‌تر از ۱ که تنها دو مقسوم‌علیه مثبت دارد،
    \emph{عدد اول} نامیده می‌شود.
    \label{def:prime}
\end{definition}

\begin{theorem}[قضیه اقلیدس]
    بی‌نهایت عدد اول وجود دارد.
    \label{thm:euclid}
\end{theorem}

\begin{lemma}[لم تقسیم‌پذیری]
    اگر $a \mid b$ و $b \mid c$، آنگاه $a \mid c$.
    \label{lem:divisibility}
\end{lemma}

\begin{corollary}[نتیجه نمونه]
    اگر $p$ عدد اول باشد، $\sqrt{p}$ گنگ است.
    \label{cor:sqrt-prime}
\end{corollary}

\begin{proposition}[گزاره نمونه]
    اگر $n$ زوج باشد، $n^2$ نیز زوج است.
    \label{prop:even-square}
\end{proposition}

%----------------------------------------
\section{مثال، نکته، تمرین}
\label{sec:example-remark-exercise}

\begin{example}[اعداد اول کوچک]
    اعداد $2, 3, 5, 7, 11, 13$ نمونه‌هایی از اعداد اول هستند.
    \label{ex:small-primes}
\end{example}

\begin{remark}[نکته]
    عدد ۱ نه اول است و نه مرکب.
    \label{rem:one-not-prime}
\end{remark}

\begin{exercise}[تمرین نمونه]
    ثابت کنید که مجموع دو عدد زوج، زوج است.
    \label{exe:sum-even}
\end{exercise}

%----------------------------------------
\section{راه حل و اثبات}
\label{sec:solution-proof}

\begin{proof}
    فرض کنید $p_1, p_2, \ldots, p_n$ همه‌ی اعداد اول باشند.
    عدد $N = p_1 p_2 \cdots p_n + 1$ را در نظر بگیرید.
    این عدد بر هیچ‌کدام از $p_i$ بخش‌پذیر نیست، پس باید
    عامل اول جدیدی داشته باشد. تناقض.
\end{proof}

\begin{solution}
    فرض کنید $a = 2m$ و $b = 2n$. آنگاه
    $a + b = 2(m + n)$ که زوج است.
\end{solution}

%----------------------------------------
\section{ارجاعات cleveref}
\label{sec:references}

ارجاعات به محیط‌های مختلف:

\begin{itemize}
    \item تعریف: \cref{def:prime}
    \item قضیه: \cref{thm:euclid}
    \item لم: \cref{lem:divisibility}
    \item نتیجه: \cref{cor:sqrt-prime}
    \item گزاره: \cref{prop:even-square}
    \item مثال: \cref{ex:small-primes}
    \item نکته: \cref{rem:one-not-prime}
    \item تمرین: \cref{exe:sum-even}
\end{itemize}

ارجاع ترکیبی: \cref{thm:euclid,lem:divisibility,cor:sqrt-prime}

%----------------------------------------
\section{شمارنده مشترک}
\label{sec:shared-counter}

\textbf{بررسی:} قضیه، لم، نتیجه، گزاره، تعریف، مثال، نکته،
تمرین **همه شمارنده‌ی مشترک** دارند. پس اگر تعریف \ref{def:prime}
شماره‌ی ۱ باشد، قضیه \ref{thm:euclid} باید شماره‌ی ۲ داشته باشد.

%----------------------------------------
\section{نتیجه‌گیری}
\label{sec:conclusion}

اگر این سند بدون خطا کامپایل شود:

\begin{itemize}
    \item ✅ محیط‌های قضیه، لم، نتیجه، گزاره کار می‌کنند
    \item ✅ تعریف، مثال، نکته، تمرین کار می‌کنند
    \item ✅ راه حل بدون شماره است
    \item ✅ اثبات با \texttt{\textbackslash qed} نمایش داده می‌شود
    \item ✅ شمارنده مشترک بین قضیه‌ها فعال است
    \item ✅ \texttt{cleveref} با نام‌های فارسی کار می‌کند
    \item ✅ رنگ‌ها (آبی، نارنجی، سبز) اعمال می‌شوند
\end{itemize}

\textbf{وضعیت تست:} قبول

\end{document}
